\section{Evaluation Protocol}
\label{sec:protocol}

Every number in this paper was produced by us on one machine
($2{\times}$RTX~A6000). All four comparison systems were built from source
locally; \emph{no number is copied from a paper}. This section states the
protocol before the results, because we will show that the protocol is worth
more than most of the differences it is used to measure.

\subsection{One protocol for every system}
Rendering quality is scored identically for all systems:
\begin{itemize}\itemsep2pt
\item \textbf{Scoring views:} $100$ held-out frames per scene, selected to be
      keyframes of \emph{no} participating system, so every scored view is
      genuinely unseen by all of them.
\item \textbf{Rendering pose:} ground-truth pose, Sim(3)-aligned~\cite{umeyama1991}
      into each map's own (arbitrary, monocular) frame.
\item \textbf{Renderer and metrics:} one code path, one PSNR and one
      LPIPS~\cite{zhang2018lpips} implementation, one exposure-normalisation
      setting, for every system.
\item \textbf{Sensor parity:} monocular against monocular.
\end{itemize}
Trajectory accuracy is ATE~RMSE via \texttt{evo\_ape tum -as}~\cite{grupp2017evo}.

\paragraph{Datasets.} Rendering is evaluated on all eight
Replica~\cite{straub2019replica} scenes and on TUM
RGB-D~\cite{sturm2012tum} \texttt{freiburg1}; trajectory accuracy on all nine
TUM \texttt{freiburg1} sequences. Head-only adaptation
(\cref{sec:exp:head}) additionally uses EuRoC~\cite{burri2016euroc},
ETH3D~\cite{schops2017eth3d} and 7-Scenes~\cite{shotton2013scenecoord},
giving five dataset families in total.

\paragraph{Alignment.}
Each monocular map lives in its own arbitrary $\mathrm{Sim}(3)$ frame, so
ground-truth poses must be transported into it. We fit
$(s,R,t)$ by Umeyama's closed-form solution~\cite{umeyama1991} on the
estimated/ground-truth camera-centre pairs and invert it,
\begin{equation}
R_{\text{map}} = R^{\top} R_{\text{gt}},\qquad
t_{\text{map}} = R^{\top}\bigl(t_{\text{gt}}-t\bigr)/s ,
\label{eq:align}
\end{equation}
noting that the $1/s$ belongs to the translation only --- folding it into the
rotation block yields a non-orthonormal matrix that silently corrupts every
rendered viewpoint. Systems that log both estimated and ground-truth keyframe
poses (MonoGS) are aligned from those pairs directly; the rest from their
trajectory files.

\subsection{A mandatory self-consistency check}
\label{sec:protocol:selfcheck}
Before any cross-system number is reported we require that \emph{a map must
render well from the poses it was built at}. This is not a formality. Our first
MonoGS adapter assumed its logged trajectory was world-to-camera when it is
camera-to-world; the inverted convention scored MonoGS at $4.36$\dB{} against
$18.71$\dB{} for the correct one, and would have produced a large, flattering,
publishable-looking margin. We adopted the rule that a result which favours us
receives one more verification step than one that does not; it changed our
conclusion three separate times.

\subsection{The protocol offset}
\label{sec:protocol:offset}

Two independent measurements quantify how much the protocol is worth on
\emph{unchanged} systems.

\noindent\textbf{(i) Same system, our protocol versus the literature.}
Photo-SLAM built from its own source scores $22.23$\dB{} on Replica office0
under our protocol; third-party comparison tables report ${\sim}30.9$\dB{} for
Photo-SLAM on monocular Replica. Offset: $\mathbf{8.7}$\dB{}.

\noindent\textbf{(ii) Same map, two pose sources.}
Taking MonoGS's own output map and changing only the rendering pose:
$18.71$\dB{} from its own estimated poses, $9.87$\dB{} from ground-truth poses.
Offset: $\mathbf{8.8}$\dB{}.

Two different systems, two different methods, essentially the same figure. We
conclude that GS-SLAM rendering numbers in the literature and ours differ by
approximately one protocol's worth of evaluation choices --- scoring on
training versus held-out views, estimated versus ground-truth poses,
resolution, undistortion, and post-sequence optimisation budget --- rather than
by quality. Concretely, our Replica $23.0$\dB{} mean should \emph{not} be read
as trailing a published $30.9$\dB{}: measured identically, we are ahead of the
system that published it.

\paragraph{What the offset is made of.}
We can attribute part of it directly. Scoring the \emph{same} map from
estimated versus ground-truth poses is worth $8.8$\dB{} on MonoGS and
$2.6$\dB{} on Photo-SLAM, the difference between the two reflecting how much
each system's map has co-adapted to its own pose error. Beyond pose source,
the protocols we found differ in: whether keyframes are excluded from scoring
(MonoGS excludes its own; Photo-SLAM's self-reported PSNR is computed
\emph{on} its keyframes), resolution ($640\!\times\!480$ undistorted versus our
$512\!\times\!384$), masking (PSNR over $I_{gt}>0$ pixels only), and
post-sequence optimisation budget ($26{,}000$ iterations for MonoGS versus our
$300$\,s polish). We do not claim to have decomposed the offset exactly; we
claim that it is large, reproducible, and larger than most reported
inter-method differences.

\paragraph{Reporting rule.} We therefore tabulate only numbers we produced
under one protocol. Published numbers are cited in text where relevant but are
never placed in a comparison table beside our own. We recommend this practice
generally for this literature.

\subsection{Systems compared}
\Cref{tab:systems} lists what each system can contribute. Two of the four
produce no renderable map and therefore appear only in the trajectory table;
MonoGS ships RGB-D-only Replica configurations, so including it there against
our monocular system would give it a depth-input advantage and we omit it
rather than report an unfair win.

\begin{table}[t]
\centering
\small
\setlength{\tabcolsep}{4pt}
\begin{tabular}{@{}llcl@{}}
\toprule
System & Type & Rend. & Role here \\
\midrule
Ours & mono 3DGS & \checkmark & --- \\
Photo-SLAM & mono 3DGS & \checkmark & rendering \\
MonoGS & mono 3DGS & \checkmark & rendering (TUM) \\
MASt3R-SLAM & mono dense & --- & tracking control \\
VGGT-SLAM & pose+cloud & --- & tracking control \\
\bottomrule
\end{tabular}
\caption{Systems compared: Photo-SLAM~\cite{huang2024photoslam},
MonoGS~\cite{matsuki2024monogs}, MASt3R-SLAM~\cite{murai2025mast3rslam},
VGGT-SLAM~\cite{maggio2025vggtslam}. The latter two produce no renderable map
--- we verified by source inspection that the VGGT-SLAM repository contains no
rasteriser or photometric metric --- so they appear only in the trajectory
table. MASt3R-SLAM is the system we inherit tracking from, making it a control
rather than an opponent.}
\label{tab:systems}
\end{table}
